\chapter{Rabies Antibody Test}

\section*{What Is This Test?}
Rabies antibody testing measures neutralizing antibodies against rabies virus. It is used mainly to verify an immune response after pre-exposure or post-exposure vaccination and in selected occupational or public-health settings.

\section*{Alternative Names}
Rabies Neutralizing Antibody; Rabies Titer; Rabies Antibody Titer; Rapid Fluorescent Focus Inhibition Test (RFFIT)

\section*{How It Is Performed}
Serum is tested by a virus-neutralization method such as RFFIT or a comparable validated assay. The result is reported as an antibody concentration or neutralization titer, commonly in international units per millilitre.

\section*{Why It Is Ordered}
\begin{itemize}
  \item Confirming an adequate response after rabies pre-exposure or post-exposure vaccination
  \item Monitoring people with ongoing occupational exposure, such as laboratory workers, veterinarians, and animal handlers
  \item Assessing suspected inadequate vaccine response in an immunocompromised person
\end{itemize}

\section*{Interpretation and Limitations}
The laboratory and public-health authority determine the protective threshold and timing of repeat testing. Antibody testing does not diagnose rabies infection and should never delay urgent post-exposure prophylaxis after a potentially rabid-animal exposure. Immunoglobulin administration, recent vaccination, immunosuppression, and timing of collection affect interpretation.
\section*{Specimen and Preparation}
Testing uses serum and does not require fasting. Interpret the collection date against the vaccination schedule and the type of exposure or occupational risk. Provide vaccine products and dates, any rabies immunoglobulin received, immune-suppressing medicines, and the public-health or occupational program requesting the test.

\section*{Risks and Safety}
Physical risk is limited to venipuncture. This is an immunity-assessment test, not a diagnostic test for clinical rabies. After a possible exposure, immediately wash the wound and contact emergency or public-health services; do not wait for an antibody result before starting recommended post-exposure care.

\section*{Understanding Results and Follow-up}
The report should give the neutralizing antibody result, units, assay method, and program cutoff. A result below the threshold may indicate inadequate or waning response, mistimed collection, or impaired immune response; it does not diagnose rabies. Follow-up, including any booster or repeat titre, must follow the responsible public-health or occupational-health protocol.

\section*{Regulatory Status and Authorization Date}
Regulatory status applies to the specific assay, instrument, manufacturer, and intended use rather than to the test name alone. No single approval date can be assigned to this test category. For a commercial assay, verify the current product in the relevant regulator's database: the U.S. Food and Drug Administration (FDA) clearance, approval, or authorization; India's Central Drugs Standard Control Organisation (CDSCO) licence or permission; the European Union In Vitro Diagnostic Regulation (EU IVDR) CE certificate issued through the applicable conformity-assessment route; or the United Kingdom Medicines and Healthcare products Regulatory Agency (MHRA) registration and UKCA/CE status. Laboratory-developed tests may not have an individual FDA, CDSCO, EU IVDR, or MHRA approval date.
\clearpage
